\documentclass[11pt]{article}
\usepackage[margin=1.1in]{geometry}
\usepackage{amsmath,amssymb,amsthm}
\usepackage[colorlinks=true,linkcolor=blue,citecolor=blue,urlcolor=blue]{hyperref}
\usepackage{booktabs}

\newtheorem{theorem}{Theorem}
\newtheorem{lemma}{Lemma}
\newtheorem{proposition}{Proposition}
\newtheorem{corollary}{Corollary}
\theoremstyle{remark}
\newtheorem*{remark}{Remark}

\newcommand{\Dfour}{\ensuremath{\mathsf{D4}}}
\newcommand{\UNK}{\ensuremath{\mathsf{UNK}}}
\newcommand{\FAL}{\ensuremath{\mathsf{FAL}}}
\newcommand{\TRU}{\ensuremath{\mathsf{TRU}}}
\newcommand{\CON}{\ensuremath{\mathsf{CON}}}
\newcommand{\meetk}{\ensuremath{\wedge_k}}
\newcommand{\joink}{\ensuremath{\vee_k}}
\newcommand{\meetR}{\ensuremath{\mathrm{meet}_R}}
\newcommand{\joinR}{\ensuremath{\mathrm{join}_R}}
\newcommand{\Clo}{\ensuremath{\mathrm{Clo}}}
\newcommand{\Eqv}{\mathrm{Eqv}}
\newcommand{\mcr}{\ensuremath{\mathrm{mcr}}}
\newcommand{\dec}{\ensuremath{\mathrm{dec}}}
\newcommand{\desc}{\mathrm{desc}}
\newcommand{\asc}{\mathrm{asc}}
\newcommand{\bits}{\mathit{bits}}

\title{The Price of NOT on \Dfour}
\author{Won Chul Yang\\ \small Seoul, Republic of Korea \\ \small \texttt{lamb5228@snu.ac.kr}}
\date{Public edition v1.0 --- 2026}

\begin{document}
\maketitle

\begin{center}\small\itshape
Revised from the 2026 preprint: the exact-cost appendix is simplified to a
single interval-shell regime via the new unconditional lemma $W_f \subseteq
B_f$ (Theorem~\ref{thm:WB}), discovered during the machine-verified audit of
the original proof. All headline results are machine-verified in Lean~4
(kernel axioms \texttt{propext}, \texttt{Quot.sound} at most; no
\texttt{native\_decide}, no \texttt{sorry}, no mathlib); see
Section~\ref{sec:mech}.
\end{center}

\begin{abstract}
This paper studies the Price of NOT on the resolved face of the four-state
dual-rail carrier \Dfour. Its sharp quantitative result is the exact
resolved-face theorem $\nu(f)=\dec(f)$. The surrounding carrier material is
used only as a compact ambient support package: the closed core
$G=\Clo\{\meetk,\joink,P\}$, the confinement-breaking involution $\sigma_s$,
the corrected $P$-equivariant-collapse theorem, and the absorbed-face
transfer needed to place the exact theorem inside one compatible carrier
presentation. Appendix~D contains the proof architecture for the exact-cost
theorem --- in this edition a single interval-shell regime, made
unconditional by the new lemma $W_f\subseteq B_f$.
\end{abstract}

\noindent\textbf{Keywords:} negation complexity; inversion complexity;
dual-rail carriers; resolved-face transfer; epistemic logic.\\
\textbf{MSC 2020:} 03B50, 03C05, 06A06, 68Q65.

\section{Introduction}\label{sec:intro}

This paper starts from the earlier resolved-block theorem of Paper~I
\cite{wcy1} and asks whether its exact price of negation persists when
viewed through a compatible four-state dual-rail carrier frame \Dfour. It
does: on the resolved block $R=\{\FAL,\TRU\}$, the minimal syntactic
negation count $\nu(f)$ is exactly the chain decrease number $\dec(f)$ for
every nonconstant $f:R^n\to R$. The larger carrier frame is admitted only
insofar as it sharpens the boundary around this exact law.

For Boolean functions, the price of negation is classical: Markov
\cite{markov} determined the circuit inversion complexity as
$\lceil\log_2(d(f)+1)\rceil$, and Morizumi \cite{morizumi} proved the
formula-level law is exact, $I_{\mathrm{for}}(f)=d(f)$. The present paper
proves the formula-exact law in the four-valued dual-rail setting, with a
proof intrinsic to the recovered resolved-face language; the relation to the
Boolean transport route is discussed at Theorem~\ref{thm:exact}.

The supporting package has six layers: (1) the \Dfour{} carrier and its
two-coordinate presentation; (2) the qualitative closed core
$G=\Clo\{\meetk,\joink,P\}$; (3) the confinement-breaking operation
$\sigma_s$; (4) the corrected expressive-collapse theorem; (5) the
gate-absorbed resolved-face transfer into Theorem~\ref{thm:exact}; (6) the
unary geometric and narrow execution specializations.

The paper does not claim that adjoining $\sigma_s$ yields all functions on
\Dfour; the correct target class is the full $P$-equivariant clone. It also
does not promote $NCP$, the Arrow of Computation, variable-cost RG, or
empirical bridge results to headline theorems. The outer
crown/constant-barrier picture belongs more naturally to a broader carrier
companion line; here it remains background, while $\nu(f)=\dec(f)$ remains
the sharp quantitative face of the paper.

The closed core preserves the internal symmetry and endpoint discipline of
the carrier, while $\sigma_s$ breaks that discipline without crossing the
rail-swap fence. On the resolved face, the earlier theorem surface is
recovered not as a plain restriction $G|_R$, but through a fixed
absorbed-face presentation compatible with the decompressed carrier theory.
This supplies an ambient gateway into the exact theorem rather than a
replacement theorem spine. The appendix structure mirrors that proof
architecture: Appendix~A records the closed-core background, Appendix~B
proves the corrected collapse, Appendix~C fixes the absorbed-face transfer,
Appendix~D contains the exact-cost proof, and Appendix~E isolates the chain
bookkeeping used there.

\paragraph{What changed in this edition.}
The audit of the original exact-cost appendix (a full machine replay of
every layer) revealed that its residual regime --- the local witness
reduction, star-local fiber coherence, and five-case rerouting layers --- is
vacuous: the containment $W_f\subseteq B_f$ holds \emph{unconditionally}
(Theorem~\ref{thm:WB}), so the interval-shell regime always applies and the
residual layers can never be reached. This edition therefore proves the
upper bound through the single interval-shell regime. The original
compressed classification arguments are superseded, not corrected: no error
was found in any reachable layer.

\section{The \Dfour{} Carrier and the Closed Monotone Core}

\subsection{Carrier}
Let $\Dfour=\{\UNK,\FAL,\TRU,\CON\}$, identified with the dual-rail carrier
$\{0,1\}^2$ by
\[
\UNK=(0,0),\quad \FAL=(0,1),\quad \TRU=(1,0),\quad \CON=(1,1).
\]
Use the coordinate functions $s=(e^++e^-)/2$ and $d=(e^+-e^-)/2$, so that
$s(\UNK)=0$, $s(\FAL)=s(\TRU)=1/2$, $s(\CON)=1$, and $d(\UNK)=0$,
$d(\FAL)=-1/2$, $d(\TRU)=1/2$, $d(\CON)=0$.

Let $P$ denote rail-swap: $P(\UNK)=\UNK$, $P(\FAL)=\TRU$, $P(\TRU)=\FAL$,
$P(\CON)=\CON$. The internal lattice operations are coordinatewise meet and
join in the knowledge order:
\[
(a,b)\meetk(c,d)=(a\wedge c,\,b\wedge d),\qquad
(a,b)\joink(c,d)=(a\vee c,\,b\vee d).
\]

\subsection{Closed Core}
Define the closed carrier core $G=\Clo\{\meetk,\joink,P\}$. In the main
text we use only the following qualitative theorem.

\begin{theorem}[Qualitative closed-core soundness]\label{thm:core}
Every term operation in $G$ is monotone for the knowledge order, preserves
both endpoints $\UNK$ and $\CON$, and is $P$-equivariant.
\end{theorem}

This is the forward-confinement background for the rest of the paper.
Stronger converse characterizations are left to Appendices~A and~B.

\section{The Confinement-Breaking Operation}

Define the simultaneous rail-complement involution
$\sigma_s(a,b)=(1-a,\,1-b)$; equivalently $\sigma_s(\UNK)=\CON$,
$\sigma_s(\FAL)=\TRU$, $\sigma_s(\TRU)=\FAL$, $\sigma_s(\CON)=\UNK$.

\begin{theorem}[Confinement-breaking algebraic move]\label{thm:sigma}
The operation $\sigma_s$ is an involution, commutes with $P$, fails to
preserve both endpoints, and is non-monotone for the knowledge order.
\end{theorem}

These are the only carrier-level facts needed later. No uniqueness claim is
used in the main theorem surface; when a stronger phrase is convenient, the
safe one is that $\sigma_s$ is a non-monotone completing primitive.

\section{Carrier Collapse with the Correct Symmetry Fence}

For each arity $n\ge 1$, let $C_n$ be the $n$-ary term clone generated by
projections together with $\meetk,\joink,P,\sigma_s$. Let $P$ act
coordinatewise on $\Dfour^n$, and define
\[
\Eqv_P(n)=\{f:\Dfour^n\to\Dfour \mid f(Px)=P(f(x))\ \text{for all}\ x\}.
\]

\begin{theorem}[Corrected expressive-collapse theorem]\label{thm:collapse}
For every arity $n\ge 1$, $C_n=\Eqv_P(n)$.
\end{theorem}

Write $x_i=(p_i,q_i)$ and $\bits(x)=(p_1,q_1,\dots,p_n,q_n)$, and let $S$
denote rail-swap on bit tuples.

\begin{lemma}[Rail-swap soundness]\label{lem:soundness}
For every $n$-ary term $t$ built from projections, $\meetk$, $\joink$, $P$,
and $\sigma_s$, there exists a Boolean function
$\phi_t:\{0,1\}^{2n}\to\{0,1\}$ such that
$t(x)=(\phi_t(\bits(x)),\,\phi_t(S(\bits(x))))$ for all $x\in\Dfour^n$. In
particular, every term operation in $C_n$ is $P$-equivariant.
\end{lemma}

\begin{lemma}[Formula lifting]\label{lem:lifting}
Let $\phi$ be a Boolean formula in the variables $p_1,q_1,\dots,p_n,q_n$.
Then there exists an $n$-ary carrier term $T_\phi$ over
$\meetk,\joink,P,\sigma_s$ such that
$T_\phi(x)=(\phi(\bits(x)),\,\phi(S(\bits(x))))$ for all $x\in\Dfour^n$.
\end{lemma}

Theorem~\ref{thm:collapse} follows by combining Lemma~\ref{lem:soundness}
with Lemma~\ref{lem:lifting}: the generators enforce exactly
$P$-equivariance, and every $P$-equivariant carrier map is recovered from a
single Boolean coordinate formula; the full proof is given in Appendix~B.

Two boundary remarks will be used later. Adjoining $\sigma_s$ yields the
maximal class compatible with rail-swap symmetry, namely the full
$P$-equivariant clone, not all carrier maps; compare the finite-set clone
viewpoint in \cite{lau,post}. On the resolved face $R=\{\FAL,\TRU\}$, the
zero-constant outer fragment is the $P$-equivariant crown fragment, and
adjoining one resolved constant completes it. This gives a coarse outer
invariant $\nu_{\mathrm{carrier}}$, but not the paper's main quantitative
theorem.

\section{Resolved-Face Transfer and Exact Cost}\label{sec:transfer}

Let $R=\{\FAL,\TRU\}$. The resolved-face theory is not obtained by plain
restriction of the internal carrier generators.

\begin{proposition}[Plain restriction is insufficient]\label{prop:plain}
On mixed resolved inputs, the internal carrier operations leave the
resolved face:
\[
\meetk(\FAL,\TRU)=\UNK,\quad \meetk(\TRU,\FAL)=\UNK,\quad
\joink(\FAL,\TRU)=\CON,\quad \joink(\TRU,\FAL)=\CON.
\]
Hence the resolved-block operations of the narrow manuscript cannot be
obtained by plain restriction of the internal carrier generators.
\end{proposition}

Fix designated re-entry data $\rho_{\mathrm{meet}}(\UNK)=\FAL$,
$\rho_{\mathrm{join}}(\CON)=\TRU$, with
$\rho_{\mathrm{meet}}|_R=\rho_{\mathrm{join}}|_R=\mathrm{id}_R$, and define
\[
\meetR(a,b)=\rho_{\mathrm{meet}}(a\meetk b),\qquad
\joinR(a,b)=\rho_{\mathrm{join}}(a\joink b).
\]

\begin{proposition}[Gate-absorbed recovery of the resolved face]
After fixing the designated re-entry data, the resolved-face operations are
recovered from carrier evaluation by absorbed re-entry on the only escaping
mixed cases.
\end{proposition}

\paragraph{Transfer principle.}
The earlier resolved-block language is not obtained by plain restriction of
$\meetk,\joink,P$ to $R$; rather, after fixing the designated re-entry data
for the two escaping mixed cases, it is recovered through a specific
absorbed-face presentation compatible with the decompressed carrier theory.
We use this presentation only as the gateway from carrier background to the
exact resolved-face theorem. The underlying table lemmas are recorded in
Appendix~C.

The exact-cost proof is phrased in order-theoretic terms: $\mcr(f)$ is the
maximal chain-reversal count and $\nu(f)$ is the resolved-face negation
complexity in the recovered term language. An alternative Boolean transport
route may be compared with Morizumi's formula-level inversion result
\cite{morizumi}, but the proof used here is intrinsic to the recovered
resolved-face language.

\begin{theorem}[Exact negation complexity on the resolved face]
\label{thm:exact}
For every nonconstant function $f:R^n\to R$, the minimal number $\nu(f)$ of
syntactic negation occurrences is exactly the chain decrease number
$\dec(f)$.
\end{theorem}

\begin{proof}[Proof outline]
For the lower bound, projections contribute no chain reversal, $\meetR$ and
$\joinR$ are reversal-subadditive, and one occurrence of $\sigma$
contributes at most one new reversal along any chain; see Appendix~D.1.
Hence $\mcr(f)\le\nu(f)$.

For the upper bound, Appendix~D supplies a one-$\sigma$ lowering step in a
\emph{single} regime: by Theorem~\ref{thm:WB} the containment
$W_f\subseteq B_f$ holds unconditionally, so the $A_f$-ideal closure of
$W_f$ is always an interval shell $S$ with $W_f\subseteq S\subseteq B_f$,
realized by one occurrence of $\sigma$, and the lowered residual satisfies
$\mcr(f_S)=\mcr(f)-1$ (Theorem~\ref{thm:decrement}). The inductive
completion in Appendix~D.3 then gives $\nu(f)\le\mcr(f)$, hence
$\nu(f)=\dec(f)$.
\end{proof}

This remains the sharp observable of the paper; the larger carrier story is
used only as compact ambient support for this theorem and should not be
read as displacing it. Appendices~B and~C prepare the recovered resolved
face, Appendix~D proves the exact theorem, and Appendix~E isolates the
auxiliary chain language used in that proof.

\section{Unary Geometry of Negation}

On the resolved face $R=\{\FAL,\TRU\}$, write $\sigma=\sigma_s|_R$. Since
both $P$ and $\sigma_s$ exchange $\FAL$ and $\TRU$, their restrictions to
$R$ agree.

\begin{proposition}[Unary resolved-face term classification]
The unary resolved-face term functions are exactly
\[
\mathrm{id}_R,\qquad \sigma,\qquad
c_{\FAL}(x)=\meetR(x,\sigma(x)),\qquad
c_{\TRU}(x)=\joinR(x,\sigma(x)).
\]
\end{proposition}

\begin{theorem}[Unary flip-localization]
In the unary dual-rail resolved slice, the restricted negation
$\sigma=\sigma_s|_R=P|_R$ is the unique flip-sign reverser.
\end{theorem}

This is the cleanest geometric specialization of the exact-cost theorem: in
the unary resolved slice, $\sigma$ is the only genuine direction-reversing
map.

\section{Executed Reversals Under Finite Episode Energy}

The execution layer remains narrow. It specializes the finite-event budget
argument of Paper~I \cite{wcy1} and does not claim a thermodynamic theorem
in the sense of the classical Landauer--Bennett discussion
\cite{landauer,bennett}. We count only executed unary resolved reversals,
namely events moving between $\FAL$ and $\TRU$ on the resolved slice.

\begin{theorem}[Finite executed reversal bound]
Under finite episode energy $E_\tau<\infty$ and a fixed positive cost
$\epsilon_{\mathrm{rev}}>0$ per executed unary flip-sign reversal, only
finitely many such executed reversals can occur, bounded by
$\lfloor E_\tau/\epsilon_{\mathrm{rev}}\rfloor$.
\end{theorem}

\begin{corollary}[Bounded reversal reachability]
A finite episode may fail to realize some unary resolved reversals that
become realizable only in a larger episode with a larger available budget.
\end{corollary}

This section does not claim a variable-cost RG theorem, a three-channel
energy theorem, a general syntax-versus-execution lower bound, or a
higher-arity execution classification.

\section{Scope and Outlook}

The paper proves a broader theorem package than the earlier resolved-block
manuscript, but the exact observable remains $\nu(f)=\dec(f)$ on the
resolved face. Carrier-level additions are admitted only if they sharpen
that observable. For the same reason, the broader carrier line is companion
background rather than the main ownership claim of the present paper. The
following remain outside the main theorem surface: $NCP$ as a headline
theorem; the Arrow of Computation as a headline theorem; variable-cost RG
or three-channel energy theoremization; empirical bridge material.

\subsection{Appendix Order}
The appendix order is: Appendix~A, qualitative and enumerative proof
material for $G$; Appendix~B, the full proof of corrected
Theorem~\ref{thm:collapse}; Appendix~C, the decompression and transfer
lemmas placing the resolved face correctly; Appendix~D, the
exact-complexity proof package for Theorem~\ref{thm:exact} (single
interval-shell regime in this edition); Appendix~E, chain-reversal, mcr,
and $\Lambda^{up}$ tools where they clarify the exact theorem.

\section{Mechanization}\label{sec:mech}

Every headline statement of this paper is machine-verified in
dependency-minimal Lean~4 (core only; no mathlib; no
\texttt{native\_decide}; no \texttt{sorry}; kernel axiom profile at most
\texttt{[propext, Quot.sound]}, with the finite table facts axiom-free).
The artifact accompanies this paper and comprises:
\begin{itemize}
\item \textbf{Theorem~\ref{thm:core}} --- \texttt{cterm\_monotone},
  \texttt{cterm\_endpoints}, \texttt{cterm\_equivariant} (structural
  induction over a closed-core term syntax);
\item \textbf{Theorem~\ref{thm:sigma}} --- \texttt{sigd\_invol},
  \texttt{sigd\_Pd}, \texttt{sigd\_breaks\_endpoints},
  \texttt{sigma\_not\_monotone} (axiom-free);
\item \textbf{Propositions C.1/C.2/C.4} --- \texttt{escape},
  \texttt{recovery}, \texttt{reentry\_independent} (axiom-free);
\item \textbf{Appendix~D lower bound} --- the per-chain subadditivity
  lemmas \texttt{chainDrops\_and}, \texttt{chainDrops\_or}, and the
  end-value-refined alternation law \texttt{chainDrops\_not}, assembled
  into \texttt{dec\_le\_sigCount};
\item \textbf{Appendix~D upper bound} --- the nested one-$\sigma$ peel is
  built on a mechanized tower library: the interval-shell decrement, the
  explicit uniform witness, and the constructive monotone extension
  (\texttt{tower\_of\_decPts}), together with the monotone-DNF realization
  (\texttt{mono\_realize}) and the tower-to-term assembly
  (\texttt{tower\_realize});
\item \textbf{Theorem~\ref{thm:exact}} --- \texttt{nu\_eq\_dec\_D4}: the
  exact equality, stated with terms evaluated on the resolved face by the
  absorbed operations $\meetR$, $\joinR$ and $\sigma$, for nonconstant $f$
  on the $n$-cube.
\end{itemize}

In addition, the original appendix was audited by exhaustive machine
replay: the closed-core and collapse layers, the transfer tables, the
interval-shell theorem (all 464 valid shells at $n\le 3$; the canonical
shell at $n=4$), and the inductive completion ($65{,}368/65{,}368$
nonconstant cases at $n\le 4$). The audit also verified computationally
that the residual of the interval regime is empty for all $n\le 4$ --- the
observation that led to Theorem~\ref{thm:WB}. The replay scripts and logs
are included in the artifact repository.

\appendix

\section{Qualitative and Enumerative Core Material}

This appendix records the forward qualitative facts about the closed core
$G=\Clo\{\meetk,\joink,P\}$ used in the main text. Their broader natural
home is the companion carrier line; here they are retained only to keep the
transfer-to-exact package self-contained.

\subsection{Proof of the qualitative soundness theorem}

\begin{proof}[Proof of Theorem~\ref{thm:core}]
Each generator is monotone for the knowledge order. The operations $\meetk$
and $\joink$ are coordinatewise meet and join on $\{0,1\}^2$, hence are
monotone. The involution $P$ is order-preserving because it only swaps the
two coordinates.

Each generator also preserves the endpoints $\UNK=(0,0)$ and $\CON=(1,1)$:
$\meetk(\UNK,\UNK)=\UNK$, $\joink(\UNK,\UNK)=\UNK$,
$\meetk(\CON,\CON)=\CON$, $\joink(\CON,\CON)=\CON$, and $P(\UNK)=\UNK$,
$P(\CON)=\CON$.

Finally, every generator is $P$-equivariant. For the involution itself this
is immediate. For $\meetk$ and $\joink$, coordinatewise meet and join
commute with swapping the two coordinates. Since monotonicity, endpoint
preservation, and $P$-equivariance are all stable under composition, every
term operation in $G$ has the same three properties.
\end{proof}

\subsection{Small carrier checks}

The proof above uses only the generator-level identities. For convenience,
we record the two endpoint rows explicitly:
\[
\UNK\meetk x=\UNK,\quad \UNK\joink x=x,\qquad
\CON\meetk x=x,\quad \CON\joink x=\CON.
\]
Thus $\UNK$ and $\CON$ behave as the lower and upper endpoints of the
internal knowledge order, while $P$ fixes those endpoints and exchanges the
resolved states $\FAL$ and $\TRU$.

\begin{remark}
The appendix does not attempt a converse description of the full core $G$.
The main text uses only the forward confinement theorem, and the stronger
carrier-collapse statement appears only after adjoining $\sigma_s$ in
Appendix~B.
\end{remark}

\subsection{Orbit detectors on the resolved face}

Work now on the resolved face $R=\{\FAL,\TRU\}$, $E=\{\UNK,\CON\}$. Every
$P$-orbit in $R^n$ has size 2, and each orbit has a unique representative
whose first coordinate is $\FAL$.

\begin{proposition}\label{prop:A1}
For $i\ne j$, define
\begin{align*}
Eq_{ij}(x)&=\joink\bigl(\meetk(x_i,x_j),\,\meetk(P(x_i),P(x_j))\bigr),\\
Neq_{ij}(x)&=\joink\bigl(\meetk(x_i,P(x_j)),\,\meetk(P(x_i),x_j)\bigr).
\end{align*}
On $R^n$, both maps are $E$-valued. Moreover, $Eq_{ij}(x)=\CON\iff
x_i=x_j$ and $Neq_{ij}(x)=\CON\iff x_i\ne x_j$.
\end{proposition}

\begin{proof}
Check the four resolved input pairs. If $x_i=x_j$, then exactly one of the
two inner meets in the definition of $Eq_{ij}$ contributes a resolved value
and the other contributes its $P$-partner, so the outer join is $\CON$;
otherwise both inner meets are $\UNK$. The same calculation with one rail
swapped gives the statement for $Neq_{ij}$.
\end{proof}

For $\alpha\in\{0,1\}^{n-1}$, define
$D_\alpha(x)=\bigwedge_{j=2}^{n}B_j(x)$, where $B_j=Eq_{1j}$ if
$\alpha_{j-1}=0$ and $B_j=Neq_{1j}$ if $\alpha_{j-1}=1$. For
$A\subseteq\{0,1\}^{n-1}$, set $D_A=\bigvee_{\alpha\in A}D_\alpha$.

\begin{proposition}
Each $D_\alpha$ is $E$-valued and satisfies $D_\alpha(x)=\CON$ exactly on
the $P$-orbit indexed by $\alpha$, while $D_\alpha(x)=\UNK$ on all other
orbits.
\end{proposition}

\begin{proof}
By Proposition~\ref{prop:A1}, each factor $B_j$ detects whether the $j$-th
coordinate agrees or disagrees with the first coordinate exactly as
prescribed by $\alpha$. Since the factors are $E$-valued, their
knowledge-meet acts as conjunction on these orbit conditions. Hence
$D_\alpha=\CON$ exactly on the orbit with pattern $\alpha$.
\end{proof}

\subsection{The outer resolved-face crown}

\begin{remark}
The broader carrier companion line proves the following resolved-face crown
fact: before adjoining any resolved constant, the $R$-valued restrictions
of $G$ on $R^n$ are exactly the $P$-equivariant maps $h:R^n\to R$, and
their count is $2^{2^{n-1}}$. In the present paper we use only the
qualitative boundary consequence that the zero-constant outer fragment
remains inside the $P$-equivariant crown.
\end{remark}

\subsection{One resolved constant completes the outer fragment}

\begin{remark}
The same companion line also proves that adjoining one resolved constant
completes the outer resolved-face fragment to all maps $R^n\to R$. The
present paper uses only the resulting outer zero/one accessibility
contrast, not the full theorem-proof package itself. This binary outer
boundary is kept only as context around the finer exact theorem
$\nu(f)=\dec(f)$.
\end{remark}

\section{Proof of the Corrected Expressive-Collapse Theorem}

This appendix expands the proof of Theorem~\ref{thm:collapse} only to the
extent needed for the present paper's self-contained
carrier-to-resolved-face gateway.

\subsection{Rail-swap normal form}

\begin{proof}[Proof of Lemma~\ref{lem:soundness}]
Proceed by induction on the term $t$. For a projection $t(x)=x_i=(p_i,q_i)$,
take $\phi_t(\bits(x))=p_i$; then
$t(x)=(\phi_t(\bits(x)),\phi_t(S(\bits(x))))$.

If the claim holds for $u$ and $v$, it is preserved by $\meetk$ and
$\joink$ because these are coordinatewise Boolean conjunction and
disjunction on $\{0,1\}^2$. Hence the first coordinate of $u\meetk v$ is
$\phi_u\wedge\phi_v$, and the first coordinate of $u\joink v$ is
$\phi_u\vee\phi_v$.

If the claim holds for $u$, then it also holds for $P(u)$ because applying
$P$ interchanges the two coordinates:
$P(u)(x)=(\phi_u(S(\bits(x))),\phi_u(\bits(x)))$. Thus one may take
$\phi_{P(u)}=\phi_u\circ S$.

If the claim holds for $u$, then it also holds for $\sigma_s(u)$ because
$\sigma_s(a,b)=(1-a,1-b)$:
$\sigma_s(u)(x)=(1-\phi_u(\bits(x)),\,1-\phi_u(S(\bits(x))))$. Hence one
may take $\phi_{\sigma_s(u)}=\neg\phi_u$.

This completes the induction. In particular, every term operation obtained
from the generators is $P$-equivariant.
\end{proof}

\subsection{Formula lifting}

\begin{proof}[Proof of Lemma~\ref{lem:lifting}]
Translate Boolean formulas recursively into carrier terms. For a positive
rail variable $p_i$, use the carrier variable $x_i$ itself. For a negative
rail variable $q_i$, use $P(x_i)$, whose first coordinate is $q_i$. For
conjunction and disjunction, use $\meetk$ and $\joink$. For negation, use
$\sigma_s$. More precisely, define a term $T_\phi$ by recursion on $\phi$:
\[
T_{p_i}=x_i,\quad T_{q_i}=P(x_i),\quad
T_{\phi\wedge\psi}=T_\phi\meetk T_\psi,\quad
T_{\phi\vee\psi}=T_\phi\joink T_\psi,\quad
T_{\neg\phi}=\sigma_s(T_\phi).
\]
An induction on the Boolean formula now shows that the first coordinate of
$T_\phi(x)$ is exactly $\phi(\bits(x))$, while the second coordinate is the
same formula evaluated on the swapped bit tuple $S(\bits(x))$.
\end{proof}

\subsection{Collapse to the full \texorpdfstring{$P$}{P}-equivariant clone}

\begin{proof}[Proof of Theorem~\ref{thm:collapse}]
The inclusion $C_n\subseteq\Eqv_P(n)$ is exactly the final sentence of
Lemma~\ref{lem:soundness}. For the converse, let $f\in\Eqv_P(n)$. Define
the Boolean coordinate function $\phi_f(\bits(x))=\pi_1(f(x))$, where
$\pi_1$ denotes the first rail coordinate. Since $f$ is $P$-equivariant,
the second rail coordinate is determined by the same function on the
swapped tuple: $f(x)=(\phi_f(\bits(x)),\,\phi_f(S(\bits(x))))$.
Lemma~\ref{lem:lifting} therefore produces a carrier term $T_{\phi_f}$ with
exactly this value, so $f\in C_n$.
\end{proof}

\section{Decompression and Transfer Lemmas}

This appendix records the table-level lemmas behind the absorbed recovered
resolved-face operations.

\subsection{Escaping mixed cases}

\begin{proposition}\label{prop:C1}
For $a,b\in R=\{\FAL,\TRU\}$, the internal carrier operations stay in $R$
except on the mixed inputs, where they escape in exactly one way:
$\meetk(\FAL,\TRU)=\meetk(\TRU,\FAL)=\UNK$ and
$\joink(\FAL,\TRU)=\joink(\TRU,\FAL)=\CON$.
\end{proposition}

\begin{proof}
This is an immediate coordinatewise calculation in the dual-rail model
$\FAL=(0,1)$, $\TRU=(1,0)$. Taking coordinatewise meet produces
$(0,0)=\UNK$, while taking coordinatewise join produces $(1,1)=\CON$. On
the unmixed inputs, one obtains the original resolved state again.
\end{proof}

\subsection{Recovered resolved-face operations}

\begin{proposition}\label{prop:C2}
With the designated re-entry data $\rho_{\mathrm{meet}}(\UNK)=\FAL$,
$\rho_{\mathrm{join}}(\CON)=\TRU$, and
$\rho_{\mathrm{meet}}|_R=\rho_{\mathrm{join}}|_R=\mathrm{id}_R$, the
absorbed recovered operations satisfy $\meetR(a,b)=\min(a,b)$ and
$\joinR(a,b)=\max(a,b)$ for all $a,b\in R$.
\end{proposition}

\begin{proof}
On unmixed pairs, Proposition~\ref{prop:C1} shows that the internal
operation already lands in $R$, so the re-entry maps act trivially. On the
two mixed pairs, the meet case lands in $\UNK$ and is sent back to $\FAL$,
while the join case lands in $\CON$ and is sent back to $\TRU$. This is
exactly the two-point chain minimum and maximum on $R=\{\FAL,\TRU\}$.
\end{proof}

\begin{corollary}
The resolved-face term language used in Section~\ref{sec:transfer} is not
the plain restriction of the internal carrier generators. It is the
absorbed $R$-face obtained by evaluating $\meetk$ and $\joink$ in the
carrier and then applying the designated re-entry data on the two escaping
mixed cases.
\end{corollary}

\begin{proof}
The failure of plain restriction is Proposition~\ref{prop:plain}. The
recovered operations are Proposition~\ref{prop:C2}.
\end{proof}

\subsection{Witness-independence at the resolved endpoints}

\begin{proposition}\label{prop:C4}
Suppose $\rho'_{\mathrm{meet}}$ and $\rho'_{\mathrm{join}}$ are alternative
designated re-entry maps with the same endpoint values:
$\rho'_{\mathrm{meet}}|_R=\mathrm{id}_R$,
$\rho'_{\mathrm{meet}}(\UNK)=\FAL$,
$\rho'_{\mathrm{join}}|_R=\mathrm{id}_R$,
$\rho'_{\mathrm{join}}(\CON)=\TRU$. Then the induced binary operations on
$R$ agree with $\meetR$ and $\joinR$ on all inputs from $R^2$.
\end{proposition}

\begin{proof}
Proposition~\ref{prop:C1} shows that the only possible outputs of $\meetk$
on $R^2$ are $\FAL,\TRU,\UNK$, and the only possible outputs of $\joink$ on
$R^2$ are $\FAL,\TRU,\CON$. The two re-entry systems agree on exactly these
values. Hence they induce identical binary operations on the resolved face.
\end{proof}

\begin{remark}
This extensional agreement does not imply uniqueness of the surrounding
carrier realization. Different decompressed witness systems may still
behave differently away from the endpoint values used here. The paper does
not count or classify that ambient witness dependence.
\end{remark}

\subsection{Transport to the recovered resolved-face language}

\begin{proposition}
Let $\sigma=P|_R$. Then the term language generated on $R$ by $\meetR$,
$\joinR$, $\sigma$ is exactly the recovered resolved-face language used in
Section~\ref{sec:transfer}.
\end{proposition}

\begin{proof}
Proposition~\ref{prop:C2} identifies $\meetR$ and $\joinR$ with the
two-point minimum and maximum on $R=\{\FAL,\TRU\}$. The unary map
$\sigma=P|_R$ exchanges the two resolved endpoints. These are precisely the
operations used in the recovered resolved-face term algebra.
\end{proof}

\begin{remark}[Safe transfer principle]
The appendix proves only the narrow statement needed in the main text: the
resolved-face language is not obtained by plain restriction, but by
absorbed re-entry on the two escaping mixed cases. It does not prove a
quotient theorem, a congruence-image theorem, or any theorem counting
ambient witness use.
\end{remark}

\section{Exact Negation Complexity on the Resolved Face}\label{app:D}

This appendix gives the proof of Theorem~\ref{thm:exact}. In this edition
the exact-cost argument has three layers: lower-bound charging, the
interval-shell decrement made unconditional by the new lemma
$W_f\subseteq B_f$, and inductive completion.

\paragraph{Comparison with the original preprint.}
The original appendix carried three further layers --- local witness
reduction, star-local fiber coherence, and a five-case rerouting package
--- to handle a residual family of shells assumed to fall outside the
interval regime. Theorem~\ref{thm:WB} below shows that this residual family
is empty: the containment $W_f\subseteq B_f$ holds for every nonconstant
$f$, so the interval-shell regime already covers every inductive step. The
removed layers were never reachable; no error in them was implicated. This
simplification was found during the machine-verified audit of the appendix.

\subsection{Lower-Bound Charging}

Let $R=\{\FAL,\TRU\}$ with the chain order $\FAL<\TRU$, and let $R^n$ carry
the componentwise order. For a maximal chain $C=(c_0<\dots<c_n)$ define
$\desc(f,C)$ to be the number of indices $i$ such that $f(c_i)=\TRU$,
$f(c_{i+1})=\FAL$. Set $\mcr(f)=\max_C \desc(f,C)$. On the two-point chain
one has $\mcr(f)=\dec(f)$.

\begin{lemma}\label{lem:D1}
If $t$ is a projection, then $\desc(t,C)=0$ for every maximal chain $C$.
\end{lemma}

\begin{proof}
A projection is monotone on $R^n$, so along any increasing chain it can
change only from $\FAL$ to $\TRU$, never from $\TRU$ to $\FAL$.
\end{proof}

\begin{lemma}\label{lem:D2}
For resolved-face terms $u$ and $v$, and every maximal chain $C$,
\[
\desc(\meetR(u,v),C)\le\desc(u,C)+\desc(v,C),\qquad
\desc(\joinR(u,v),C)\le\desc(u,C)+\desc(v,C).
\]
\end{lemma}

\begin{proof}
If $\meetR(u,v)$ descends at a step of $C$, then both inputs are $\TRU$
below and at least one is $\FAL$ above, so that descent is charged to an
input descent. If $\joinR(u,v)$ descends, then both inputs are $\FAL$ above
and at least one is $\TRU$ below, so again the descent is charged to an
input descent. Summing over all steps gives the inequalities.
\end{proof}

\begin{lemma}\label{lem:D3}
For every resolved-face term $u$ and maximal chain $C$,
$\desc(\sigma(u),C)\le\desc(u,C)+1$.
\end{lemma}

\begin{proof}
Along a fixed chain, applying $\sigma$ exchanges $\FAL$ and $\TRU$, so
descents of $\sigma(u)$ are ascents of $u$. In any binary word, the number
of ascents exceeds the number of descents by at most one.
\end{proof}

\begin{proposition}[Lower-bound charging]\label{prop:D4}
If a resolved-face term $t$ computes $f:R^n\to R$, then $\mcr(f)\le\nu(t)$.
In particular, $\dec(f)=\mcr(f)\le\nu(f)$.
\end{proposition}

\begin{proof}
Proceed by structural induction on $t$. Lemma~\ref{lem:D1} handles
projections. Lemma~\ref{lem:D2} handles $\meetR$ and $\joinR$.
Lemma~\ref{lem:D3} handles one occurrence of $\sigma$, which can add at
most one new descent along any chain. Therefore any term computing $f$ has
at least $\mcr(f)$ occurrences of $\sigma$. Taking the minimum over all
such terms gives $\mcr(f)\le\nu(f)$, and $\mcr(f)=\dec(f)$ on the two-point
chain.
\end{proof}

\subsection{Interval-Shell Decrement}

Assume now that $f:R^n\to R$ is nonconstant and non-monotone. Define
\[
A_f=\{x: f(x)=\TRU\},\qquad
G_f={\uparrow}A_f\setminus A_f,\qquad
B_f=A_f\setminus{\uparrow}G_f.
\]
Let $W_f$ be the union of the first $\TRU$-blocks over all worst chains of
$f$. For $S\subseteq R^n$, define the lowered residual $f_S$ by
$f_S(x)=\FAL$ for $x\in S$, $f_S(x)=f(x)$ otherwise.

\begin{proposition}[One-$\sigma$ realization for interval shells]
\label{prop:D5}
Let $S$ be an $A_f$-ideal with $W_f\subseteq S\subseteq B_f$. Then $\chi_S$
is definable with at most one occurrence of $\sigma$.
\end{proposition}

\begin{proof}
Set $U_S={\uparrow}S$, $E_S=U_S\setminus S$. Then
$S=U_S\setminus{\uparrow}(E_S)$. If $E_S=\varnothing$, then $S=U_S$ is an
upset and $\chi_S$ is negation-free. Otherwise,
$\chi_S=\meetR(\chi_{U_S},\,\sigma(\chi_{{\uparrow}(E_S)}))$, so one
occurrence of $\sigma$ suffices.
\end{proof}

\begin{lemma}[Worst-chain shell identity]\label{lem:D6}
If $C$ is a worst chain for $f$, then $B_f\cap C$ is exactly the first
$\TRU$-block of $f$ on $C$.
\end{lemma}

\begin{proof}
Let the first descent on $C$ occur at $c_j\to c_{j+1}$. Then
$c_{j+1}\in G_f$, so every point strictly above the first $\TRU$-block lies
in ${\uparrow}G_f$ and hence outside $B_f$. Points below that block are
outside $A_f$, so they are also outside $B_f$. The first $\TRU$-block
itself lies in $B_f$.
\end{proof}

\begin{lemma}[Chain-prefix property]\label{lem:D7}
Let $S$ be an $A_f$-ideal with $S\subseteq B_f$. Then on every maximal
chain $C$, the set $S\cap C$ is a lower prefix of each $\TRU$-block of $f$
on $C$.
\end{lemma}

\begin{proof}
If $x<y$ lie in the same $\TRU$-block on $C$ and $y\in S$, then $y\in B_f$.
Lemma~\ref{lem:D6} implies that $x$ also lies in $B_f$. Since $x\in A_f$
and $S$ is downward closed inside $A_f$, it follows that $x\in S$.
\end{proof}

\begin{corollary}[No new descents]\label{cor:D8}
Under the hypotheses of Lemma~\ref{lem:D7}, for every maximal chain $C$,
$\desc(f_S,C)\le\desc(f,C)$.
\end{corollary}

\begin{proof}
Lowering only shortens lower prefixes of existing $\TRU$-blocks, so it
cannot create a new $\TRU\to\FAL$ boundary.
\end{proof}

\begin{theorem}[Exact one-step decrement for interval shells]
\label{thm:decrement}
Let $S$ be an $A_f$-ideal with $W_f\subseteq S\subseteq B_f$. Then
$\mcr(f_S)=\mcr(f)-1$.
\end{theorem}

\begin{proof}
Fix a worst chain $C$ for $f$. Since $W_f\subseteq S$, the first
$\TRU$-block on $C$ lies in $S$. Since $S\subseteq B_f$,
Lemma~\ref{lem:D6} forces $S\cap C$ to lie inside that same first block.
Hence $S\cap C$ is exactly the first $\TRU$-block, and deleting it lowers
the descent count on $C$ by exactly one.

For any other maximal chain $D$, Corollary~\ref{cor:D8} gives
$\desc(f_S,D)\le\desc(f,D)$. If $D$ is non-worst, then
$\desc(f,D)\le\mcr(f)-1$. If $D$ is worst, the same first-block argument
gives $\desc(f_S,D)=\mcr(f)-1$. Thus all chains have at most $\mcr(f)-1$
descents after lowering, and at least one worst chain has exactly that
many.
\end{proof}

\begin{corollary}[Forward interval-shell safety]\label{cor:D10}
Every $A_f$-ideal $S$ with $W_f\subseteq S\subseteq B_f$ gives a one-step
safe shell: it is realized by one occurrence of $\sigma$, and lowering
along it decreases $\mcr$ by exactly one.
\end{corollary}

\begin{proof}
Combine Proposition~\ref{prop:D5} with Theorem~\ref{thm:decrement}.
\end{proof}

\begin{theorem}[$\bigstar$ Unconditional containment $W_f\subseteq B_f$]
\label{thm:WB}
For every nonconstant, non-monotone $f:R^n\to R$, $W_f\subseteq B_f$.
\end{theorem}

\begin{proof}
Suppose not: some $w\in W_f$ lies outside $B_f$. Since
$w\in W_f\subseteq A_f$, this means $w\in{\uparrow}G_f$, so there is a gap
point $g\in G_f$ with $g\le w$, and by the definition of $G_f$ a point
$a\in A_f$ with $a\le g$, $f(a)=\TRU$, $f(g)=\FAL$. Both containments are
strict: $a\ne g$ because $f(a)\ne f(g)$, and $g\ne w$ because $f(g)=\FAL$
while $f(w)=\TRU$. So $a<g<w$.

Since $w\in W_f$, there is a worst chain $C$ for $f$ whose first
$\TRU$-block contains $w$. On $C$, no descent occurs before the first
$\TRU$-block, so all $\mcr(f)$ descents of $C$ occur on the segment of $C$
from $w$ upward.

Build an increasing chain $X$: ascend to $a$, then to $g$, then to $w$,
then follow $C$ upward from $w$. The segment through $a<g<w$ contributes at
least one descent (its $f$-word passes from $\TRU$ at $a$ to $\FAL$ at
$g$), and the tail from $w$ contributes $\mcr(f)$ descents. By the
refinement lemma (Lemma~\ref{lem:E4}), some maximal chain has at least as
many descents as $X$, so $\mcr(f)\ge\desc(f,X)\ge 1+\mcr(f)$, a
contradiction.
\end{proof}

\begin{corollary}[Canonical shell --- the interval regime is total]
\label{cor:D12}
For every nonconstant, non-monotone $f$, the $A_f$-ideal closure of $W_f$,
$S_f=A_f\cap{\downarrow}W_f$, satisfies $W_f\subseteq S_f\subseteq B_f$.
Hence the interval-shell decrement of Corollary~\ref{cor:D10} applies at
every inductive step: every nonconstant non-monotone $f$ admits a
one-$\sigma$ shell whose lowering decreases $\mcr$ by exactly one.
\end{corollary}

\begin{proof}
$S_f$ is an $A_f$-ideal containing $W_f$ by construction. For
$S_f\subseteq B_f$, let $x\in S_f$, so $x\in A_f$ and $x\le w$ for some
$w\in W_f$. If $x$ were outside $B_f$, there would be a gap point
$g\in G_f$ with $g\le x\le w$, contradicting $w\in B_f$
(Theorem~\ref{thm:WB}).
\end{proof}

\subsection{Inductive Completion}

\begin{proposition}[Monotone base case]\label{prop:D13}
Let $f:R^n\to R$ be nonconstant. If $\mcr(f)=0$, then $f$ is monotone and
therefore has a negation-free resolved-face term.
\end{proposition}

\begin{proof}
If $\mcr(f)=0$, then no increasing chain contains a descent
$\TRU\to\FAL$, so $f$ is monotone. Hence $A_f$ is an upset. Since $f$ is
nonconstant, the usual monotone formula
$f=\bigvee_{a\in\mathrm{Min}(A_f)}\ \bigwedge_{i:\,a_i=\TRU}x_i$ computes
$f$ using only $\meetR$ and $\joinR$.
\end{proof}

\begin{theorem}[Upper bound]\label{thm:D14}
For every nonconstant function $f:R^n\to R$, $\nu(f)\le\mcr(f)$.
\end{theorem}

\begin{proof}
We induct on $\mcr(f)$. If $\mcr(f)=0$, Proposition~\ref{prop:D13} gives
$\nu(f)=0$.

Assume $\mcr(f)>0$ and the claim already holds for all nonconstant
functions of smaller $\mcr$. Since $\mcr(f)>0$, $f$ is non-monotone, so by
Corollary~\ref{cor:D12} the canonical shell $S_f$ is an $A_f$-ideal with
$W_f\subseteq S_f\subseteq B_f$, and Corollary~\ref{cor:D10} gives a
one-$\sigma$ lowering step whose residual $f_1=f_{S_f}$ satisfies
$\mcr(f_1)=\mcr(f)-1$. One occurrence of $\sigma$ realizes the first step
and reconstruction uses only monotone context. Hence
$\nu(f)\le 1+\nu(f_1)$. By the induction hypothesis,
$\nu(f_1)\le\mcr(f_1)=\mcr(f)-1$. Therefore $\nu(f)\le\mcr(f)$.
\end{proof}

\begin{corollary}[Exact equality]
For every nonconstant $f:R^n\to R$, $\nu(f)=\mcr(f)=\dec(f)$.
\end{corollary}

\begin{proof}
The lower bound is Proposition~\ref{prop:D4} and the upper bound is
Theorem~\ref{thm:D14}. Since $\mcr(f)=\dec(f)$ on the two-point chain, all
three quantities are equal.
\end{proof}

\section{Auxiliary Chain-Reversal Tools}

This appendix records a small amount of chain notation used as bookkeeping
for Appendix~\ref{app:D}.

\subsection{Upward edges and descending edges}

Let $X=(x^{(0)}<x^{(1)}<\dots<x^{(m)})$ be an increasing chain in $R^n$.
Write $\Lambda^{up}(X)=\{(x^{(j)},x^{(j+1)}): 0\le j<m\}$ for its set of
adjacent upward edges. For $f:R^n\to R$, define the descending edge set
$\Delta_f(X)=\{(u,v)\in\Lambda^{up}(X): f(u)>f(v)\}$. Then
$\mcr_X(f):=|\Delta_f(X)|$ is the chainwise reversal count, and
$\mcr(f)=\max_X \mcr_X(f)$ is the global maximum over increasing chains.

\subsection{Ascents and descents in binary chain words}

Let $w=(w_0,\dots,w_m)\in R^{m+1}$ be a binary word. Write $\asc(w)$ for
the number of indices $j$ with $w_j=\FAL$ and $w_{j+1}=\TRU$, and
$\desc(w)$ for the number of indices $j$ with $w_j=\TRU$ and
$w_{j+1}=\FAL$.

\begin{lemma}\label{lem:E1}
For every binary word $w$, $\asc(w)\le\desc(w)+1$.
\end{lemma}

\begin{proof}
Read the word from left to right and decompose it into maximal constant
blocks. Every ascent starts a $\TRU$-block, and every descent ends one. At
most one $\TRU$-block can fail to be preceded by a descent, namely the
initial $\TRU$-block if the word begins at $\TRU$. Hence the number of
ascents exceeds the number of descents by at most one.
\end{proof}

\begin{corollary}
Let $X$ be an increasing chain and let $\sigma$ exchange $\FAL$ and $\TRU$.
Then $\mcr_X(\sigma\circ f)\le\mcr_X(f)+1$.
\end{corollary}

\begin{proof}
Along the chain $X$, descents of $\sigma\circ f$ are ascents of the binary
word defined by $f$. Apply Lemma~\ref{lem:E1}.
\end{proof}

\subsection{Prefix lowering lemma}

\begin{lemma}\label{lem:E3}
Fix a binary word along an increasing chain, and lower some $\TRU$-entries
to $\FAL$. If the lowered set is a lower prefix inside each maximal
$\TRU$-block, then no new descending edge is created.
\end{lemma}

\begin{proof}
Inside a single maximal $\TRU$-block, lowering a lower prefix changes a
block of the form $\TRU\cdots\TRU$ into
$\FAL\cdots\FAL\,\TRU\cdots\TRU$. This creates at most one new ascent
inside the block and never a new descent. At the block boundaries, the
left boundary can only move from $\FAL\,\TRU$ to $\FAL\,\FAL$, and the
right boundary can only move from $\TRU\,\FAL$ to either $\TRU\,\FAL$ or
$\FAL\,\FAL$. Hence the number of descents does not increase.
\end{proof}

\subsection{Subchain refinement lemma}

\begin{lemma}\label{lem:E4}
Every increasing chain $X$ in $R^n$ extends to a maximal chain $C$ with
$\desc(f,C)\ge\mcr_X(f)$. Hence the maximum over increasing chains in the
definition of $\mcr$ agrees with the maximum over maximal chains.
\end{lemma}

\begin{proof}
Refine each upward edge $u<v$ of $X$ into a saturated cover path from $u$
to $v$ and concatenate, extending below the bottom of $X$ and above its top
to a maximal chain $C$. If $(u,v)$ is a descending edge of $X$, the
$f$-word along the inserted cover path starts at $\TRU$ and ends at $\FAL$,
so it contains at least one descending cover step. Distinct edges of $X$
refine to disjoint segments of $C$, so
$\desc(f,C)\ge|\Delta_f(X)|=\mcr_X(f)$.
\end{proof}

\begin{remark}
Appendix~\ref{app:D} uses Lemmas~\ref{lem:E3} and~\ref{lem:E4} in the
interval-shell argument and in Theorem~\ref{thm:WB}. They are recorded
separately here only to keep the main exact-cost appendix shorter.
\end{remark}

\begin{thebibliography}{9}
\bibitem{cohomnote} W.~C. Yang. \emph{The cohomological price of NOT.}
Zenodo, 2026. DOI: 10.5281/zenodo.21775055. (Companion note; cites the
present paper and mechanizes the nested tower used in Appendix~D's Lean
artifact.)
\bibitem{markov} A.~A. Markov. \emph{On the inversion complexity of a
system of functions.} J.~ACM 5(4), 331--334, 1958. (Russian original:
Doklady AN SSSR 116(6), 917--919, 1957.)
\bibitem{morizumi} H.~Morizumi. \emph{Limiting Negations in Formulas.}
ICALP 2009, LNCS 5555, 701--712. (Preliminary version: arXiv:0811.0699,
2008.)
\bibitem{bennett} C.~H. Bennett. \emph{Logical reversibility of
computation.} IBM J.~Res.~Dev. 17(6):525--532, 1973.
\bibitem{landauer} R.~Landauer. \emph{Irreversibility and heat generation
in the computing process.} IBM J.~Res.~Dev. 5(3):183--191, 1961.
\bibitem{lau} D.~Lau. \emph{Function Algebras on Finite Sets.} Springer,
2006.
\bibitem{post} E.~L. Post. \emph{The Two-Valued Iterative Systems of
Mathematical Logic.} Princeton University Press, 1941.
\bibitem{wcy1} W.~C. Yang. \emph{Finite-Energy Epistemic Logic with
Conservative Pointed Extension and Negation Geometry.} 2026. (Paper~I;
companion release: \url{https://github.com/ycmath/finite-energy-epistemic-logic}.
An earlier draft circulated as ``Finite-energy epistemic logic with
T0-preserving open updates''.)
\end{thebibliography}

\section*{Authorship \& provenance}

Won Chul Yang, independent researcher. This paper and its companion Lean
artifact were produced in collaboration with an AI research loop operated
by the author (frontier language models --- Anthropic Claude family --- for
discovery, formalization, and adversarial verification), \textbf{with the
Lean~4 kernel as the final acceptance gate}. The interval single-regime
simplification of Appendix~D (Theorem~\ref{thm:WB}) was found during a
machine audit of the original proof. The author directed the research
programme and verified the pipeline; in line with the author's
research-ethics policy, no claim of academic priority is made beyond full
disclosure of how the work was produced. Corrections are invited.

Licenses: Apache-2.0 (Lean artifacts and verification scripts), CC~BY~4.0
(text).

\end{document}
